% ==========================================
% 第 4 章：学习路线
% ==========================================
\section{学习路线与开发基础}

\begin{quotebox}
\small
学习机器人开发时，不必一开始就深入所有底层细节。对于 ROS 2、Linux 与 Git 等工具，应当先理解它们要解决的工程问题，再依据所承担的实际开发任务逐步深入。\\
既不能完全依赖生成式工具输出未经推导的代码，也不应陷入漫无目的地逐行阅读依赖库的低效怪圈。最有效的起点，是先亲手跑通一条端到端的完整功能链路，再围绕遇到的具体故障定向深挖。\\
禁止完全许愿式编程，但也不必转向一行行了解代码，这就像从一个低效走向另一个低效。\\
\hfill —— andinm\_hong
\end{quotebox}

新成员可根据当前确定的研发方向选择阅读路径：

\begin{table}[H]
\centering
\small
\begin{tabularx}{\textwidth}{p{2.2cm} X p{4.2cm}}
\toprule
\textbf{研发方向} & \textbf{必读技术章节} & \textbf{可以先跳过的部分} \\
\midrule
感知组 & 4.4 Linux、4.5 YOLO 目标检测、4.7 粒子滤波定位，结合 3.3 IMU 与 3.4 里程计 & 4.8 深度强化学习数学推导 \\
\addlinespace
策略组 & 4.3 ROS 2 通信机制、4.6 行为树，结合 1.6 裁判盒状态机 & 4.7 粒子滤波内部积分数学推导 \\
\addlinespace
运控组 & 4.4 Linux、4.6 行为树动作接口，结合 3.2 关节自由度与 3.3 IMU 姿态解算 & 4.5 视觉检测网络训练与数据增强细节 \\
\bottomrule
\end{tabularx}
\caption{按研发方向规划的学习路径。公共基础为 Git 协作、Linux 命令行与 ROS 2 通信模型。}
\end{table}

\subsection{可以暂缓深入的内容}

之所以建议新人在入队初期暂缓钻研以下方向，是因为在尚未建立整机数据流全貌之前，过早陷入孤立的底层细节极易引发严重的挫败感，迫使学习进度停滞在细枝末节上：

\begin{itemize}[leftmargin=*]
    \item \textbf{DDS 底层报文协议细节。} 初期只需掌握节点自动发现、话题发布订阅与服务调用的基本使用，只有在遭遇局域网内广播风暴或跨网段断连等网络疑难杂症时，才需要抓包分析 RTPS 报文头。
    \item \textbf{Executor 调度器的内部源码。} 先掌握单线程与多线程执行器的调度差异，明确耗时计算会阻塞同一执行器内的其他回调即可，无需通读其源码调度循环。
    \item \textbf{C++ 高阶模板元编程。} 现有工程中大量运用了模板封装，阅读时能看懂接口调用形式即可，避免在语法技巧上耗费过多时间。
    \item \textbf{强化学习算法收敛性证明。} 重点理解训练流程、动作空间与观测空间的定义方式，数学证明不影响工程调试。
    \item \textbf{复杂文档排版技巧。} 集中精力于技术事实与工程逻辑的准确表达，排版工具按需查阅即可。
\end{itemize}

\subsection{Git：建立严格的工程协作习惯}

在多人协同开发大型机器人工程时，Git 不仅是版本备份工具，更是维护代码稳定性的第一道防线：

\begin{table}[H]
\centering
\small
\begin{tabularx}{\textwidth}{p{3.2cm} p{4.6cm} X}
\toprule
\textbf{任务场景} & \textbf{推荐命令} & \textbf{操作目的与注意事项} \\
\midrule
查看工作区变动 & \texttt{git status}、\texttt{git diff} & 提交前逐行比对修改，避免误将本地临时测试日志或中间产物混入暂存区 \\
\addlinespace
精细化暂存 & \texttt{git add -p} & 交互式逐块审查改动，杜绝将无关修改强行打包 \\
\addlinespace
原子化提交 & \texttt{git commit -m "..."} & 保持单次提交的语义原子性，便于后续故障二分回溯（\texttt{git bisect}） \\
\addlinespace
同步远程分支提交 & \texttt{git pull --rebase} & 维持清晰线性的提交历史，避免产生大量无意义的分支合并节点 \\
\addlinespace
开辟特性分支 & \texttt{git switch -c feat/...} & 隔离新功能开发，维持主分支始终处于可编译、可测试状态 \\
\addlinespace
丢弃工作区修改 & \texttt{git restore <file>} & 放弃未提交的临时改动；若需撤销已推送提交，必须使用 \texttt{git revert} 生成逆向提交 \\
\addlinespace
更新私有子模块 & \texttt{git submodule update ...} & 使底层依赖库与私有技术手册同步至指定版本校验码 \\
\bottomrule
\end{tabularx}
\caption{常用 Git 命令及其操作规程。}
\end{table}

团队协作中必须严格恪守三条纪律：
\begin{enumerate}[leftmargin=*]
    \item \textbf{严禁对公共受保护分支执行强行推送（\texttt{force push}）。} 强行覆盖远程分支将直接破坏其他协作者的本地提交基底，导致灾难性的代码丢失与繁琐的重组代价。
    \item \textbf{坚持单一职责原则。} 严禁在一个提交中同时混合“功能开发”与“全工程代码格式化”。这会导致代码审查完全无法聚焦，在遭遇缺陷需要回溯时，也会直接剥夺原子化撤销的可能性。
    \item \textbf{提交 PR 前必须在本地通过全套构建与静态检查。} CI 自动化工作流仅作为辅助防线，开发者必须对自身推送代码的本地编译与基本功能负责。
\end{enumerate}

\subsection{ROS 2 的常用机制与适用场景}

ROS 2 是机器人软件的神经系统。各个核心机制的设计直接对应着特定的工程解题需求：

\begin{table}[H]
\centering
\small
\begin{tabularx}{\textwidth}{p{2.6cm} p{5.0cm} X}
\toprule
\textbf{核心机制} & \textbf{解决的工程问题} & \textbf{常见调试与故障场景} \\
\midrule
DDS / RTPS & 去中心化节点自动发现与高吞吐网络传输 & 局域网内某台机器人重启后，排查节点是否能自动恢复拓扑连接 \\
\addlinespace
节点与通信模型 & 模块独立运行，通过话题、服务与动作交换数据 & 梳理系统架构与数据调用关系，划分模块职责边界 \\
\addlinespace
QoS 服务质量策略 & 规定消息可靠性、队列缓冲深度与保活要求 & 解决相机画面延迟过高、控制指令堆积或比赛状态偶发丢包 \\
\addlinespace
Executor 执行器 & 调度回调函数并发线程与优先级 & 解决控制周期抖动、高频通信丢帧与多回调死锁问题 \\
\addlinespace
Lifecycle 生命周期 & 提供明确的节点状态机（未配置/激活/失活） & 比赛上场前顺序自检与停机安全保护 \\
\addlinespace
Node Composition & 支持多节点在同一进程中运行，实现零拷贝通信 & 解决高分辨率高帧率图像在进程间拷贝占用过高 CPU 资源的问题 \\
\addlinespace
TF2 坐标变换系统 & 维护机器人机体、传感器与球场的空间树形转换 & 将图像中识别到的目标映射为球场绝对坐标 \\
\addlinespace
rosbag2 数据录制 & 完整记录现场传感器与话题通信时序数据 & 赛后离线回放现场数据复现偶发故障 \\
\bottomrule
\end{tabularx}
\caption{ROS 2 核心机制与实际调试场景对应关系。}
\end{table}

\noindent\textbf{QoS 配置策略的选择依据：}

\begin{itemize}[leftmargin=*]
    \item \textbf{视频图像流：} 相机传输 1080P/720P 原始帧时数据通量极高。必须采用 \texttt{BEST\_EFFORT}（尽力而为）传输模式并将队列深度设为 1。之所以这样设置，是因为在图像处理链路中，下一帧新图像永远比上一帧更具时效价值；如果误配置为 \texttt{RELIABLE}（可靠传输），一旦出现单帧网络丢包，通信中间件将不断发起重传重试，由此导致历史画面帧在缓冲区发生严重堆积，造成机器人看到的视野产生数秒的时延滞后。
    \item \textbf{裁判盒与状态机指令：} 必须采用 \texttt{RELIABLE} 模式。比赛状态切换指令低频但极为致命，漏收一次便会导致机器人无法按时开球或违规受罚，绝不允许丢失。
    \item \textbf{高频 IMU 姿态反馈：} 通常采用 \texttt{BEST\_EFFORT} 传输且队列设为 1，使运动控制回调函数获取到的数值始终保持当前最新。
\end{itemize}

\subsection{Linux 与终端开发基础}

Ubuntu Linux 是机器人机载主控的标准操作系统，熟练掌握命令行是执行真机联调的基础：

\begin{table}[H]
\centering
\small
\begin{tabularx}{\textwidth}{p{3.2cm} p{4.2cm} X}
\toprule
\textbf{命令分类} & \textbf{常用命令} & \textbf{实际工程应用场景} \\
\midrule
进程与资源监控 & \texttt{htop}、\texttt{ps aux} & 排查机载 CPU 占用是否过载、识别异常僵死进程 \\
\addlinespace
系统日志排查 & \texttt{journalctl -u ... -f} & 实时追踪开机自启系统服务的报错与崩溃追踪信息 \\
\addlinespace
磁盘空间排查 & \texttt{df -h}、\texttt{du -sh *} & 避免录制大量 rosbag2 数据导致工控机固态硬盘爆满死机 \\
\addlinespace
网络端口诊断 & \texttt{ip addr}、\texttt{netstat} & 确认调试网口 IP 配置、抓取裁判盒 3838 广播端口报文 \\
\addlinespace
软硬件权限管理 & \texttt{sudo usermod -aG ...} & 解决 USB 串口与 CAN 驱动板无权限读写的底层错误 \\
\bottomrule
\end{tabularx}
\caption{机载 Linux 系统维护核心命令。}
\end{table}

\subsection{YOLO 目标检测：从标注到推理}

视觉感知的核心目标是从图像像素中提取结构化特征。训练与部署 YOLO 必须遵循严谨的工程闭环：现场视频采集抽帧 $\rightarrow$ 标注足球与球门 $\rightarrow$ 数据增强与色温抖动 $\rightarrow$ 模型训练与 mAP 评估 $\rightarrow$ 导出 ONNX / TensorRT $\rightarrow$ 真机推理与时延评测。

\begin{enumerate}[leftmargin=*]
    \item \textbf{样本多样性压倒模型复杂度。} 之所以要在人工草皮现场采集多时间段样本，是因为训练数据集一旦缺乏强反光、强逆光或大面积阴影样本，即使选用超大规模的模型架构，在真实赛场上也会频繁发生漏检。
    \item \textbf{实时性是第一底线。} 视觉检测运行在机载 i7-1355U 上，算力资源十分紧张。检测耗时必须严格压缩在 30 毫秒以内（至少维持 30Hz 推理速度），否则将直接加剧控制回路的总时延。通常建议将 PyTorch 权重导出为 ONNX，并利用 TensorRT 或 OpenVINO 执行 FP16 量化加速。
\end{enumerate}

\subsection{行为树：任务规划与状态解耦}

战术策略层通常采用行为树（BehaviorTree）替代复杂的嵌套状态机：

\begin{table}[H]
\centering
\small
\begin{tabularx}{\textwidth}{p{2.4cm} p{3.2cm} X}
\toprule
\textbf{节点类型} & \textbf{典型节点} & \textbf{调度机制与工程作用} \\
\midrule
控制节点 & \texttt{Fallback} & 顺序尝试执行各子节点，只要有一个子节点返回 SUCCESS 便立刻返回，用于处理容错兜底 \\
\addlinespace
控制节点 & \texttt{Sequence（起身）} & 仅当所有子节点均返回 SUCCESS 时才返回成功，用于执行倒地起身固定序列 \\
\addlinespace
条件节点 & \texttt{IsFallen} & 仅读取 IMU 姿态判定机体是否倒地，只读查询，严禁引入改变外部状态的副作用 \\
\addlinespace
动作节点 & \texttt{GetUp} & 触发底层运动控制器执行起身轨迹 \\
\addlinespace
控制节点 & \texttt{Sequence（进攻）} & 顺序调度寻球、接近球体与射门动作 \\
\addlinespace
动作节点 & \texttt{FindBall / Kick} & 执行旋转寻球、调整身位逼近与单腿踢球动作 \\
\bottomrule
\end{tabularx}
\caption{行为树核心节点及其调度机制。}
\end{table}

编写行为树时必须遵守两项核心准则：
\begin{itemize}[leftmargin=*]
    \item \textbf{条件节点绝对严禁引入副作用。} 条件节点在行为树每个 tick 周期都会被高频轮询。如果开发者在检查“是否看到足球”的同时顺手修改了全局目标变量，当行为树发生逻辑回溯时，将引发不可预测的循环震荡。
    \item \textbf{严格匹配决策 tick 周期与底层控制频率。} 行为树通常以 10 至 20 Hz 的低频运行，而底层运动控制运行在 200 Hz。行为树下发的瞬态速度跳变必须经过平滑加速度插值器过滤，避免阶跃速度指令直接导致电机过流或机体失稳。
\end{itemize}

\subsection{粒子滤波与场地定位}

足式里程计能够推算短时位移，但打滑和累计误差会持续增大；前置相机能看到场地白线、中圈和球门，但因为足球场两头形状完全一致，容易认错半场。蒙特卡洛粒子滤波（MCL）在地图坐标系中维护一组候选位置点（粒子），用里程计推算它们怎么移动，再用看到的标线和球门给每个粒子打分，最后算出机器人在球场上的具体坐标。

\subsubsection{定位计算步骤}

\begin{enumerate}[leftmargin=*]
    \item \textbf{初始撒点：} 在球场范围内生成大量随机粒子，每个粒子记录一个候选坐标与朝向角度 $(x, y, \theta)$。
    \item \textbf{叠加运动推算：} 读取里程计位移增量，把机体走过的距离加到每个粒子上，并加入运动高斯噪声模拟打滑。
    \item \textbf{比对视觉标线：} 把相机检测到的白线交点、中圈弧线和球门立柱，与各粒子视角下应该看到的地图特征比对，重合度越高给的分数（权重）越高。
    \item \textbf{重采样筛选：} 淘汰得分过低的粒子，在得分高的粒子周围重新复制生成新粒子。
    \item \textbf{计算当前坐标：} 统计聚集在一起的高分粒子重心，输出机器人在球场坐标系中的当前位置与朝向。
\end{enumerate}

\begin{table}[H]
\centering
\small
\begin{tabularx}{\textwidth}{p{2.6cm} p{3.2cm} X}
\toprule
\textbf{调参项目} & \textbf{典型范围} & \textbf{调整影响与权衡取舍} \\
\midrule
粒子数量 & 500--2000 个 & 粒子过多拖慢机载工控机计算；粒子过少容易跟丢位置 \\
\addlinespace
运动噪声参数 & 根据实测打滑程度标定 & 设置过小会导致跟丢真实轨迹；设置过大会让粒子散得太开 \\
\addlinespace
观测噪声参数 & 根据相机标定精度设定 & 设置过小会过度轻信单次画面抖动；设置过大会忽略看到的有效标线 \\
\addlinespace
重采样阈值 & 有效粒子降到总数一半 & 重采样太频繁会丧失寻找备选位置的能力，太慢会导致坐标更新迟钝 \\
\bottomrule
\end{tabularx}
\caption{粒子滤波关键调优参数及其工程权衡。}
\end{table}

\noindent\textbf{常见定位问题与处理方法：}

\begin{itemize}[leftmargin=*]
    \item \textbf{半场认反。} 足球场关于中线完全对称，两头的白线和球门形状一模一样。如果机器人开场撒点不确定，或者被撞倒爬起后丢失了朝向，仅靠看到的白线就会把两个半场搞混，把自己的位置算到对方半场去。实际处理时需要结合机载指南针航向角锁定攻防方向，或者在开赛前根据裁判盒下发的队伍半场信息初始化朝向。
    \item \textbf{标线被遮挡。} 机器人身处开阔中场或视线被其他机器人挡住时，视野里找不到有效标线。此时依靠里程计做短时推算，暂停粒子重采样，避免在缺乏证据时把粒子胡乱收拢。
    \item \textbf{碰撞后重新定位。} 遭受猛烈冲撞导致机器人完全丢失当前位置时，重新在全场撒点，旋转机身寻找多个标线特征重新计算位置。
\end{itemize}

\subsection{深度强化学习的适用范围}

本节为进阶选读内容，梳理双足运动控制领域常见的基础概念：

\begin{table}[H]
\centering
\small
\begin{tabularx}{\textwidth}{p{3.2cm} p{4.6cm} X}
\toprule
\textbf{核心概念} & \textbf{数学物理内涵} & \textbf{在双足控制中的具体角色} \\
\midrule
贝尔曼方程 & 当前状态价值等于即时奖励加上经折减的未来期望价值 & 构成值函数迭代优化的理论基石 \\
\addlinespace
时序差分（TD） & 基于单步或多步转移样本更新当前价值函数 & 避免等待完整对抗回合结束，大幅提升样本利用效率 \\
\addlinespace
策略网络 & 以高频观测数据为输入，直接输出各关节目标动作 & 端到端映射状态到力矩，替代传统复杂的运动学解析解 \\
\addlinespace
Actor-Critic 架构 & Actor 网络生成动作，Critic 网络评估打分 & 降低策略梯度估计方差，促使训练平稳收敛 \\
\addlinespace
PPO 算法 & 严格限制每次策略更新的概率分布剪切步长 & 防止单次梯度更新幅度过大导致策略崩溃，在双足步态训练中应用最为广泛 \\
\bottomrule
\end{tabularx}
\caption{强化学习常见基础概念及其控制角色。}
\end{table}

\noindent\textbf{部署到真机前的物理制约与风险评估：}

\begin{itemize}[leftmargin=*]
    \item \textbf{真机跌落与机械磨损代价。} 双足机器人在真机上直接试错成本极高。每次高动态跌落都会对高减速比行星齿轮箱与轴承施加剧烈剪切载荷，因此策略训练与初步筛选必须在 Isaac Gym 或 MuJoCo 等高并发物理仿真引擎中完成。
    \item \textbf{仿真与现实差距（Sim2Real Gap）。} 仿真接触模型与实际人工草皮摩擦力、电机驱动器响应延迟存在差异。训练时必须引入域随机化（Domain Randomization）对质量分布、摩擦系数与通信延迟施加扰动。
    \item \textbf{黑盒策略可解释性。} 端到端神经网络难以直观分析特定动作选择的具体原因，现场调试与紧急修复难度较大。通常建议结合残差控制或基于经典控制框架进行策略融合。\advice
\end{itemize}

\newpage